\section{Half-Wave Symmetry}

\subsubsection*{Motivation and Intuition}

Many periodic engineering signals have additional symmetry beyond being even or odd. One important type is half-wave symmetry.

A waveform has half-wave symmetry if shifting it by half a period reverses its sign.

\begin{conceptbox}
Let $f(x)$ be periodic with period $2L$. The function has half-wave symmetry if
\[
f(x+L)=-f(x).
\]
\end{conceptbox}

This means the second half of the waveform is the negative of the first half.

\subsubsection*{Physical Interpretation}

Half-wave symmetry commonly occurs in AC waveforms, square waves, inverter outputs, and switching circuits. It implies that the positive and negative halves of the waveform balance each other.

\begin{noteBox}
If a signal has half-wave symmetry, its average value over one period is zero. Therefore, its DC component is absent:
\[
a_0=0.
\]
\end{noteBox}

\subsubsection*{Key Property}

\begin{conceptbox}
If $f(x)$ has half-wave symmetry, then all even harmonics vanish. That is,
\[
a_{2n}=0,\qquad b_{2n}=0.
\]
Only odd harmonics may remain.
\end{conceptbox}

\subsubsection*{Worked Example}

\begin{examplebox}
Consider the square wave
\[
f(x)=
\begin{cases}
k, & 0<x<\pi,\\
-k, & -\pi<x<0,
\end{cases}
\]
with period $2\pi$.

This function is odd, so
\[
a_0=0,\qquad a_n=0.
\]

We compute
\[
b_n=\frac{1}{\pi}\int_{-\pi}^{\pi}f(x)\sin(nx)\,dx.
\]

Since $f(x)$ is odd and $\sin(nx)$ is odd, their product is even:
\[
b_n=\frac{2}{\pi}\int_0^\pi k\sin(nx)\,dx.
\]

Thus,
\[
b_n=\frac{2k}{\pi}
\int_0^\pi \sin(nx)\,dx.
\]

Compute:
\[
b_n=
\frac{2k}{\pi}
\left[
-\frac{\cos(nx)}{n}
\right]_0^\pi.
\]

So,
\[
b_n=
\frac{2k}{n\pi}
\left[1-\cos(n\pi)\right].
\]

Since
\[
\cos(n\pi)=(-1)^n,
\]
we have
\[
b_n=
\frac{2k}{n\pi}\left[1-(-1)^n\right].
\]

If $n$ is even,
\[
b_n=0.
\]

If $n$ is odd,
\[
b_n=\frac{4k}{n\pi}.
\]

Therefore,
\[
f(x)=\frac{4k}{\pi}
\left[
\sin x+\frac{1}{3}\sin(3x)+\frac{1}{5}\sin(5x)+\cdots
\right].
\]
\end{examplebox}

\subsubsection*{Engineering Insight}

\begin{noteBox}
Half-wave symmetry is important in power electronics. If an inverter output has half-wave symmetry, then even harmonics are absent. This simplifies harmonic analysis and reduces certain types of distortion.
\end{noteBox}

\subsubsection*{Common Mistakes and Notes}

\begin{conceptbox}
Half-wave symmetry is not the same as odd symmetry.

Odd symmetry means
\[
f(-x)=-f(x).
\]

Half-wave symmetry means
\[
f(x+L)=-f(x).
\]

A function may have one, both, or neither.
\end{conceptbox}